\documentclass[11pt,a4paper]{article}
\usepackage[margin=2.2cm]{geometry}
\usepackage[T1]{fontenc}
\usepackage{lmodern}
\usepackage{booktabs}
\usepackage{tabularx}
\usepackage{enumitem}
\usepackage{amsmath}
\PassOptionsToPackage{hyphens}{url}
\usepackage[hidelinks]{hyperref}
\sloppy
\setlist{nosep}
\usepackage{microtype}
\setlength{\emergencystretch}{3em}
\newcolumntype{Y}{>{\raggedright\arraybackslash}X}
\newcommand{\f}[1]{\texttt{#1}}

\title{Design Specification\\[2pt]\large Part 1 as a study: the PhonePe Pulse state panel, 2018--2026}
\author{BTP, IIT Kharagpur --- Supervisor: Prof.\ Rudra Prakash Pradhan (VGSOM)}
\date{4 November 2026}

\begin{document}
\maketitle

\section{Request and understanding}

\paragraph{Request (verbatim).} ``youve framed the parts as progress updates. present them as a very
detailed study of that part. put in graphs numbers everything. write it in normal, detailed prose like a
paper, but keep following the simple english words usage. please rewrite part 1 keeping this in mind.''

\paragraph{What the request says.}
\begin{itemize}
  \item The Part 1 document changes from a two-page progress report to a detailed study of the
        Part 1 subject.
  \item The study has figures, tables and numbers.
  \item The text is paper prose: full paragraphs that explain the data, the method, each figure and
        table, and what the numbers show.
  \item The words and sentences follow ASD-STE100 Simplified Technical English (STE).
  \item Only Part 1 is rewritten now.
\end{itemize}

\paragraph{Assumptions (not said in the request).}
\begin{itemize}
  \item The reader is the supervisor. The study must stand alone: a reader who has not seen the code
        can follow the data, the method and the results.
  \item The subject of Part 1 is the adoption data and its denominator: the state crosswalk, the
        PhonePe Pulse ingestion and the population projections. The study describes these data and
        what they show in nominal terms.
  \item Parts 2 and 3 will be rewritten later in the same format. This design makes the format
        reusable; it does not plan those parts.
\end{itemize}

\paragraph{Success criteria.}
\begin{enumerate}
  \item The document reads as a study by the author (Devansh Gupta). It does not refer to ``the
        student'', to AI, or to plan task numbers.
  \item Every number in the text, the tables and the figures comes from committed data by committed
        code. No number in the prose is typed by hand.
  \item One command sequence rebuilds every number, table and figure of the study from the raw data:
        \f{make data analysis report}.
  \item The study stays inside the Part 1 data scope (Section~\ref{sec:scope}).
  \item The files that Parts 2 and 3 use do not change. The check is
        \f{git diff --exit-code} on \f{data/interim/phonepe.csv}, \f{data/interim/population.csv},
        \f{data/final/} and \f{output/tables/} after the rebuild. The figure PDFs of Part 3 are not
        compared: each PDF has its creation date inside the file.
  \item The document compiles with \f{latexmk} with no undefined references, no undefined numbers and
        no overfull boxes wider than 1\,pt.
\end{enumerate}

\section{Scope}\label{sec:scope}

\paragraph{In scope.}
\begin{itemize}
  \item State crosswalk: \f{src/py/btpdata/states.py}, \f{state\_crosswalk.csv} (46 name variants, 35
        \f{state\_id} values, main sample of 30).
  \item PhonePe Pulse at the pinned commit \f{257bf38} (24 September 2026): state-quarter transaction
        count, amount and registered users (\f{data/interim/phonepe.csv}), and the new split of the
        transaction count by type (Section~\ref{sec:types}).
  \item NCP 2020 population projections, Table 8 (\f{data/manual/population\_ncp2020.csv},
        \f{data/interim/population.csv}).
  \item Nominal values, per-person values with the NCP population, shares, growth rates, rank
        statistics, dispersion of logs across states, and data checks.
\end{itemize}

\paragraph{Out of scope.}
\begin{itemize}
  \item CPI deflation and real values (Part 2).
  \item The comparison with NPCI state data (Part 2). The study cites the NPCI data constraint (state
        data only from April 2025) as the reason for the proxy, with no NPCI numbers.
  \item Maps (they need the Part 2 boundary file).
  \item $\beta$-convergence regressions and the coefficient-of-variation figure (Part 3).
  \item Edits to \f{thesis/btp1/}, \f{thesis/btp1-slides/} and \f{thesis/progress/}.
  \item Parts 2 and 3 of the study.
\end{itemize}

\section{Approaches considered}

\begin{description}
  \item[A. Code-generated study (selected).] Python adds the new PhonePe input and the data checks.
        One R script reads the interim files and writes the figures (PDF), the tables (\f{.tex}) and a
        file of computed numbers (\f{numbers.tex}). The document gets every number in the prose from that
        file with a macro \verb|\val{key}|. An unknown key stops the compile.
        Advantage: numbers cannot drift from the data; the rebuild is one command.
        Cost: one more R script and a key list to keep.
  \item[B. Hand-typed numbers with command comments.] The method of the progress reports: each number is
        typed, and a \verb|%| comment gives the command. Rejected: a study of 20 or more pages has
        hundreds of numbers. Each data change needs a manual edit, and a typed number can differ from
        the data without an error.
  \item[C. knitr (\f{.Rnw}) with inline R.] Rejected: \f{latexmk} does not run knitr, so the
        \f{report} target needs a new step; analysis code and prose mix in one file; a reviewer cannot
        read the numbers without a build.
\end{description}

\section{Design}

\subsection{File layout}
\begin{tabularx}{\textwidth}{@{}l Y@{}}
\toprule
File & Purpose \\
\midrule
\f{src/py/btpdata/fetch\_phonepe.py} (modify) & add two sparse paths; apply the sparse set again if it
  differs from the checkout \\
\f{src/py/btpdata/parse\_phonepe\_types.py} (new) & transaction count by type; data checks \\
\f{data/interim/phonepe\_types.csv} (new) & \f{state\_id, period, variable, value, source};
  variables \f{phonepe\_count\_p2p}, \f{phonepe\_count\_retail}, \f{phonepe\_count\_utility} \\
\f{data/interim/phonepe\_checks.csv} (new) & one row per check: \f{check, n\_compared, n\_mismatch,
  max\_abs\_diff, note} \\
\f{src/R/lib/state\_labels.csv} (new) & \f{state\_id, name, code, in\_main}: display name, two-letter code
  and main-sample flag for the 35 ids \\
\f{src/R/lib/study\_theme.R} (new) & ggplot2 theme and PDF save helper for the study figures \\
\f{src/R/lib/part1\_fns.R} (new) & measures, national series, shares, concentration, decomposition,
  rank statistics \\
\f{src/R/lib/latex\_out.R} (new) & number formatting, \f{numbers.tex} writer, \f{.tex} table writer \\
\f{src/R/part1\_study.R} (new) & reads interim data; writes all study outputs \\
\f{output/study/part1/figures/*.pdf} (new) & study figures \\
\f{output/study/part1/tables/*.tex} (new) & study tables (tabular bodies, no float) \\
\f{output/study/part1/numbers.tex} (new) & computed numbers as \texttt{\textbackslash val} definitions \\
\f{thesis/study/part1.tex}, \f{part1.pdf} (new) & the study document and its committed PDF \\
\f{tests/py/test\_fetch\_phonepe.py}, \f{test\_phonepe\_types.py} (new) & Python tests \\
\f{tests/R/test\_part1.R}, \f{test\_latex\_out.R} (new) & R tests \\
\f{Makefile} (modify) & \f{data}: add \f{parse\_phonepe\_types}; \f{analysis}: add
  \f{part1\_study.R}; \f{report}: build \f{thesis/study/part*.tex} \\
\bottomrule
\end{tabularx}

\subsection{PhonePe inputs}\label{sec:types}
The pinned commit has nine data folders: \f{aggregated}, \f{map} and \f{top}, each for
\f{transaction}, \f{user} and \f{merchant}. The current clone has only \f{map/transaction/hover} and
\f{map/user/hover}. The study uses:
\begin{itemize}
  \item \f{map/transaction/hover/country/india/<y>/<q>.json}: count and amount by state (already used).
  \item \f{map/user/hover/country/india/<y>/<q>.json}: registered users by state (already used).
  \item \f{aggregated/transaction/country/india/state/<state>/<y>/<q>.json}: count by type
        (\f{P2P}, \f{Retail}, \f{Utility}). There is no amount by type at this commit. The README of the
        commit defines the types as person-to-person transfers, payments to merchants, and recharges and
        bill payments, and says that the three types sum to the total.
  \item \path{aggregated/transaction/country/india/<y>/<q>.json} and
        \path{aggregated/user/country/india/<y>/<q>.json}: national totals, for checks only.
\end{itemize}
The two new sparse paths are \path{data/aggregated/transaction/country/india} and
\path{data/aggregated/user/country/india}.

\paragraph{Facts verified on a scratch clone of the pinned commit (4 November 2026).}
\begin{itemize}
  \item The type files have 1224 state-quarters (36 raw names $\times$ 34 quarters). The sum of the
        three types equals the hover count in all 1224. One entry is absent: Lakshadweep, 2018Q1,
        \f{Utility}; the other two types already sum to the total.
  \item The national transaction files equal the sum of the 36 state counts in all 34 quarters. The
        national user files equal the sum of the 36 state user counts in all 34 quarters.
  \item The state folder names use hyphens (\f{jammu-\&-kashmir}).
  \item No state has a fall in registered users from one quarter to the next.
\end{itemize}

\paragraph{Parser rules.}
\begin{itemize}
  \item Folder name to \f{state\_id}: replace \f{-} by a space, then \f{to\_state\_id}. The crosswalk
        and \f{normalise} do not change.
  \item An absent type in a state-quarter gets 0 only if the present types sum to the hover count of that
        state-quarter. Otherwise the parser stops and names the file.
  \item The parser stops if the type sum differs from the hover count, or if a national file differs
        from the state sum. It writes \f{phonepe\_checks.csv} with the counts of compared items and of
        mismatches (0 after a passing run) and the number of absent entries set to 0.
  \item Merged ids use \f{sum\_merged}, as in \f{parse\_phonepe.py}.
\end{itemize}

\paragraph{Not used, with reason (stated in the study).}
\begin{itemize}
  \item Registered merchants: a life-till-date count; 18 of the 36 raw units have no file for one or
        more quarters of 2018--2019.
  \item District files (\f{map/.../state/...}) and top-10 lists: the unit of the BTP is the state; the
        district names change over the period and need a separate district crosswalk.
  \item App opens and device brands: not in the pinned commit. The commit restates the series from
        2018Q1 with new definitions; its README says that the series must not be joined with earlier
        releases.
\end{itemize}

\subsection{Measures}
For state $s$ and quarter $t$: $C$ is the transaction count, $A$ the amount in rupees (nominal), $U$ the
registered users, $P$ the population, and $C^k$ the count of type $k$.
\begin{align*}
\text{count per person} &= C/P &
\text{amount per person} &= A/P \\
\text{users per 1000 persons (reach)} &= 1000\,U/P &
\text{transactions per user (intensity)} &= C/U \\
\text{average transaction value} &= A/C &
\text{type share} &= C^k / C
\end{align*}
The identity $\ln(C/P) = \ln(U/P) + \ln(C/U)$ splits adoption into reach and intensity. Across the
states of the main sample in one quarter,
\[
1 = \frac{\operatorname{cov}(\ln(U/P), \ln(C/P))}{\operatorname{var}(\ln(C/P))}
  + \frac{\operatorname{cov}(\ln(C/U), \ln(C/P))}{\operatorname{var}(\ln(C/P))}.
\]
The two terms are the contributions of reach and intensity to the cross-state variance. A scratch run
gives reach 0.78 and intensity 0.22 in 2018Q4, and reach 0.39 and intensity 0.61 in 2026Q2. The study
reports these numbers from \f{numbers.tex}, not from this specification.

\paragraph{Samples and reference quarters.} Figures and tables show all 35 ids and mark the five
union territories outside the main sample. Statistics across states (variance, rank correlation,
correlation, decomposition) use the 30 ids of the main sample. State comparisons use 2018Q4 and
2026Q2.

\subsection{Study document: structure}
\f{thesis/study/part1.tex}; article class, 11pt, A4, the preamble of the progress reports plus
\f{tikz}, \f{longtable} and \f{caption}; \f{plainnat} with \f{references/refs.bib}. Title: ``Digital
payment adoption across Indian states, 2018--2026: building the state panel from PhonePe Pulse''.
Subtitle line: ``BTP study, Part 1''. Author: Devansh Gupta. Target length: 20 to 30 pages.

\begin{enumerate}
  \item \textbf{Abstract.} Data, method, four to six main numbers.
  \item \textbf{Introduction.} The BTP question; why a state panel of adoption is necessary; NPCI
        state data only from April 2025 (\f{npci\_upi\_state}, \f{medianama2025npci}); PhonePe Pulse as
        the source, with its prior use (\f{reddy2024icrier}); the aim and the contents of this part;
        what Parts 2 and 3 add (one sentence each).
  \item \textbf{Data sources.} PhonePe Pulse: what the repository is, licence, layout, the files that
        the study uses and does not use, with reasons (Table: data files). NCP 2020: what the report
        projects, the reference date (1 March), the years used. Figure: data flow (TikZ) from the raw
        files to the interim tables.
  \item \textbf{Construction of the panel.} State units and the crosswalk (Table: the 35 units, in the
        appendix); merged ids and the summation rule; parsing; the data checks (Table: checks, from
        \f{phonepe\_checks.csv} and the population checks); population interpolation (formula, NCP sum
        check, Figure: population growth by state); the measures (Section~\ref{sec:types} formulas).
  \item \textbf{Adoption in India as a whole.} Figure: national count, amount and users. Table: annual
        totals and growth. Figure: quarter-on-quarter growth (2020Q2 is the only quarter with a fall in
        the national count). Figure: average transaction value (national, and median and interquartile
        range across states). Figure: national type shares.
  \item \textbf{Adoption across states.} Figure: count per person by state, 2018Q4 and 2026Q2.
        Table: state values in 2026Q2. Table: concentration (top five states: shares of count, amount,
        users and population; number of states for half of all transactions). Rank persistence
        (Spearman correlation of count per person, 2018Q4 against 2026Q2). Figure: reach against
        intensity, 2026Q2. Figure and table: the variance decomposition over time. Figure: retail share
        against count per person, 2026Q2, with the Spearman coefficient.
  \item \textbf{Limitations of the data.} One app (its share of UPI is a Part 2 question); registered
        users are a cumulative count, not active users; how PhonePe assigns a transaction to a state is
        not documented in the repository; the population values are projections from the 2011 Census;
        values are nominal; the type shares change sharply in 2018.
  \item \textbf{Conclusions.} The findings, each with its number; what the later parts need from this
        panel.
  \item \textbf{Appendix A.} The 35 units (long table). \textbf{Appendix B.} Reproduction: commands,
        pinned commit, file paths, and the rule that \f{numbers.tex} is generated.
\end{enumerate}
Each figure and table has a paragraph in the text that says what it shows and what the numbers mean.
Interpretation is allowed when the numbers support it. A cause is stated only with evidence; else the
text says that the data do not show the cause.

\subsection{Numbers mechanism}
\begin{itemize}
  \item \path{output/study/part1/numbers.tex} has one line per number:
        \verb|\defval{nat.count.2026Q2}{38.66}| and a \verb|%| comment that names the R function and
        the input file. The header says: generated by \f{Rscript src/R/part1\_study.R}; do not edit.
  \item The document defines \verb|\defval| and \verb|\val|. \verb|\val{key}| prints the value;
        an unknown key gives \verb|\PackageError| with the key name, so the compile fails.
  \item Keys are lower case, dot-separated: \f{<topic>.<item>[.<period>]}. The plan lists every key.
  \item The R writer formats the values (thousands separators as \verb|{,}|, fixed decimals per key),
        stops on a duplicate key, on a key outside the pattern, and on a non-finite value.
  \item Tables are \f{.tex} files with the \f{tabular} body only. The document gives the float,
        caption and label.
\end{itemize}

\subsection{Figures}
R and ggplot2 (no new packages; \f{ggrepel} is not installed, so point labels are the two-letter
codes of \f{state\_labels.csv}). One theme in \f{study\_theme.R}, with the palette of
\f{src/R/01\_descriptives.R}. Width 6.3\,in. Each figure has a caption with the unit, the sample, the
period and the source. Files: \f{national\_totals}, \f{national\_growth}, \f{avg\_value},
\f{type\_shares}, \f{state\_count\_pc}, \f{reach\_intensity}, \f{decomposition}, \f{retail\_share},
\f{population\_growth} (all \f{.pdf}). The data-flow figure is TikZ in the document.

\subsection{Language rules}
\begin{itemize}
  \item STE vocabulary and sentence rules: common words, one meaning per word, short sentences
        (descriptive sentences up to 25 words), active voice, one topic per paragraph.
  \item Paper prose: paragraphs, not lists, for the method and the results. Lists only for sets of
        items (files, checks).
  \item No filler, no idioms, no hype words, no value judgments. A banned-word check (plan) catches
        common cases.
  \item Units: transactions in billion or per person per quarter; amounts in rupees (Rs.) or Rs.\ lakh
        crore ($10^{12}$ rupees); users in million or per 1000 persons; population in million.
\end{itemize}

\subsection{Testing}
\begin{itemize}
  \item Python (pytest): fetch re-applies a changed sparse set (local Git fixture repository, no
        network); type parser on fixtures: values, merged ids, hyphen names, zero fill, stop on a sum
        mismatch, stop on a national mismatch, checks file content.
  \item R (testthat): each function in \f{part1\_fns.R} on small hand-made data with known answers
        (the decomposition terms sum to 1; top-$k$ share; growth; rank correlation); the writer in
        \f{latex\_out.R} (format, duplicate key, bad key, non-finite value).
  \item Build: \f{make test}; \f{make data analysis}; \f{git diff --exit-code} on the Part 2--3 files;
        \f{latexmk} log check; banned-word check.
\end{itemize}

\subsection{Format for Parts 2 and 3 (not planned now)}
The same pattern applies: \f{thesis/study/part<N>.tex}, \f{src/R/part<N>\_study.R},
\path{output/study/part<N>/{figures,tables}/}, \path{output/study/part<N>/numbers.tex}, the same
\verb|\val| mechanism and theme, and a tag \f{part<N>-study}. Part 2 adds the CPI deflation, the NPCI
comparison and maps; Part 3 adds the convergence analysis. Each later part cites the earlier study parts
and does not repeat their tables.

\section{Decisions taken automatically}
The process asks the user at these points. No user was present; each decision below is the
recommended option.

\begin{enumerate}
  \item \textbf{Process path: architectural.} New ingestion, new analysis script and a new document
        format change several components.
  \item \textbf{Identity of the rewritten Part 1: new file \f{thesis/study/part1.tex} and new tag
        \f{part1-study}.} The tag \f{part1} and \f{thesis/progress/part1.tex} stay as the record of the
        first submission; a new directory keeps the old and new formats apart and gives Parts 2--3 a
        place.
  \item \textbf{Approach A (code-generated numbers).} Only this approach makes a wrong number a build
        error.
  \item \textbf{Use the transaction-type split.} It decomposes the main adoption variable exactly and
        is inside the PhonePe ingestion scope.
  \item \textbf{Do not use merchants, districts, top lists.} Incomplete early coverage (merchants);
        the unit is the state and districts need a new crosswalk (districts, top lists).
  \item \textbf{No CPI, no NPCI numbers, no maps.} They belong to Part 2 data; the request keeps
        Part 1 inside its own data.
  \item \textbf{No $\beta$-regression and no CV figure; use the reach--intensity variance
        decomposition.} The CV of count per person is already in Part 3; the decomposition is new and
        uses only Part 1 data.
  \item \textbf{Statistics across states use the 30-id main sample; displays show all 35.} Same
        sample rule as the master design.
  \item \textbf{Reference quarters 2018Q4 and 2026Q2.} The master plan uses 2018Q4 to avoid the very
        small 2018Q1 base; the type shares also move sharply in the first three quarters of 2018.
  \item \textbf{Python for ingestion, R for analysis and outputs.} Same split as the existing
        pipeline; no new packages.
  \item \textbf{Hyphen handling in the type parser only.} A change to \f{normalise} could change other
        sources and the Part 2--3 files.
  \item \textbf{Absent type entry is 0 only after the sum check.} The zero then follows from the data,
        not from an assumption.
  \item \textbf{New theme file; \f{01\_descriptives.R} unchanged.} A refactor of the Part 3 script
        would rebuild Part 3 figures for no change in content.
  \item \textbf{State codes for labels from ISO 3166-2:IN, logged as a source.} \f{ggrepel} is not
        installed; full names overlap.
  \item \textbf{Units: billion transactions, Rs.\ lakh crore, million users.} Readable totals;
        lakh crore is the unit of the earlier reports.
  \item \textbf{Committed PDF \f{thesis/study/part1.pdf}.} Same rule as the progress reports.
  \item \textbf{Tag \f{part1-study} created locally, not pushed.} The task forbids a push.
\end{enumerate}

\section{Risks}
\begin{tabularx}{\textwidth}{@{}Y Y@{}}
\toprule
Risk & Response \\
\midrule
GitHub fetch of the new sparse paths fails or is slow & \f{fetch} stops with the Git error; run again.
  The parser stops if a file is absent. \\
A check fails on the real data & The parser stops and names the file. Investigate before any change to
  the rule; record the finding in the study. \\
Figure PDFs differ at each build (creation date) & Rebuild comparison covers CSV, \f{.tex} tables and
  \f{numbers.tex}, not figure PDFs (same rule as Part 3). \\
Prose drifts from STE & Banned-word check; reviewer checks sentence length and voice. \\
\bottomrule
\end{tabularx}

\end{document}
